% Part: turing-machines 
% Chapter: undecidability
% Section: introduction

\documentclass[../../../include/open-logic-section]{subfiles}

\begin{document}

\olfileid{tur}{und}{int}
\olsection{引言}

并非每个函数——甚至并非每个算术函数——都是可计算的，这似乎显而易见。这样的函数实在太多了，其行为也太过复杂——例如，由放射性粒子衰变，或其他混沌或随机行为所定义的函数。假设我们从某个给定时刻开始以1秒为间隔进行计数，并定义函数 $f(n)$ 为在初始时刻之后第 $n$ 个1秒间隔内宇宙中衰变的粒子数。这看起来像是一个我们永远无法指望去计算的函数。

但是，无法想象如何去计算这些函数是一回事，实际证明它们是不可计算的则是另一回事。事实上，即使看起来无比复杂的函数，在抽象意义上也可能是可计算的。例如，假设宇宙在时间上是有限的——正如一些宇宙学理论所预言的那样，在遥远的未来某一天，宇宙将坍缩为一个点。那么，从那个初始时刻算起，使得 $f(n)$ 有定义的秒数就只有有限个（尽管大得难以置信）。而任何只对有限多个输入有定义的函数都是可计算的：我们可以把输出列在一张大表里，或者将其编码到一个非常大的图灵机状态转移图中。

我们常常对一类特殊的函数感兴趣，它们的值给出了是/否问题的答案。例如，问题“$n$ 是素数吗？”关联于函数
\[
\fn{isprime}(n) = \begin{cases}
  1 & \text{若 $n$ 是素数}\\
  0 & \text{否则。}
  \end{cases}
\]
我们说一个是/否问题是\emph{能行可判定的}，如果与之关联的取值为1/0的函数是能行可计算的。

为了从数学上证明存在无法能行计算的函数，或无法能行判定的问题，必须固定一个具体的计算模型，并证明存在它无法计算或判定的函数或问题。例如，我们可以证明并非所有函数都能被图灵机计算，也并非所有问题都能被图灵机判定。然后，我们可以诉诸 邱奇--图灵论题来得出结论：不仅图灵机的计算能力不足以计算每一个函数，而且没有任何能行过程能够做到这一点。

证明这种否定结果的关键在于，我们可以为图灵机本身分配编号。最简单的方法是枚举它们，例如，通过固定一种书写图灵机及其程序的特定方式，然后以系统化的方式列出它们。一旦我们认识到这可以做到，那么图灵不可计算函数的存在性就由简单的势的考虑得出：从 $\Nat$ 到~$\Nat$ 的函数集（实际上，甚至只是从 $\Nat$ 到 $\{0, 1\}$ 的函数集）是不可数的，但由于我们可以枚举所有图灵机，图灵可计算函数的集合只是可数无穷的。

我们也可以定义\emph{具体的}函数和问题，并分别证明它们是不可计算的与不可判定的。其中一个这样的问题就是所谓的\emph{停机问题}。图灵机可以通过列出其指令来进行有限描述。这样一种对图灵机的描述，即一个图灵机程序，当然可以作为另一台图灵机的输入。因此，我们可以考虑那些用来判定关于其他图灵机的问题的图灵机。一个特别有意思的问题是：“给定的图灵机在输入~$n$ 上启动后最终会停机吗？”如果存在一台能够判定此问题的图灵机，那将会很好：可以把它想象成一个质量管控图灵机，确保图灵机不会陷入无限循环之类的情况。图灵 证明的有趣事实是，不可能存在这样的图灵机。不可能存在这样一台单一的图灵机：当其以一台图灵机 $M$ 的描述和某个数~$n$ 作为输入启动时，总能根据机器 $M$ 在输入 $n$ 上启动后是否会停机，而相应地停机并输出 $1$ 或 $0$。

一旦我们有了具体的不可判定问题的例子，就可以用它们来证明其他问题也是不可判定的。例如，一个著名的不可判定问题是：“一阶公式~$!A$ 是否有效？”。不存在这样的图灵机：给定一个一阶公式~$!A$ 作为输入，它总能保证停机并根据 $!A$ 是否有效来输出 $1$ 或 $0$。历史上，寻找一种能行解决该问题的方法这一问题，被简单地称为“判定问题”；因此我们说，判定问题是不可解的。邱奇 和 图灵 大约在同一时期独立地证明了这个结果，所以它也被称为 邱奇--图灵 定理。

\end{document}